\section{Project overview}
\textbf{Project title.} Who Said So? Payee provenance for AI-agent payments on UPI.

\textbf{One-line summary.} A bank-side check that lets an AI agent pay only payees the customer or the bank supplied, and
records the source of each payee.

\textbf{Website description.} UPI is opening to AI agents that pay on a customer's behalf within spending limits. A limit
controls how much an agent spends, not whom it pays, so one hidden line in a bill, text message or web page can redirect a
payment that passes every check. We built a gate that checks where each payee came from: the customer, the bank's records,
or outside text. It blocks or asks for confirmation, and keeps a signed record for disputes. On banking traces from four AI
vendors, it stopped all 197 successful attacks without blocking a legitimate payment.

\section{The problem}
\textbf{DPI system addressed.} UPI, including bill payments through Bharat Connect. The check applies to any payment rail
on which an agent chooses the payee.

\textbf{Which DPI assumption breaks.} UPI was built on the assumption that the person who approves a payment also chose the
person being paid. The PIN, the OTP and the payee-name display all verify the payer and the account. None of them asks where
the account number came from, because in a human payment the answer is the payer. Agent payments change this. Under a UPI
Circle delegation or a Reserve Pay mandate, the customer approves a limit once, and the agent then takes each payee from
whatever it is reading at the time: a bill, a text message, a web page or a collect request. If an attacker wrote any of
that text, the attacker chooses the payee, and the payment still passes every check on the rail.

\textbf{Problem statement and who is affected.} Consider a customer who asks an agent to pay the electricity bill. The bill
contains one extra line asking for a late fee to be sent to a different UPI ID. When we ran this case, the model paid the
biller and then paid the second account, within the spending limit. Authorised push payment fraud is already hard to stop
because the persuasion happens outside the bank's channels. Agents add a channel in which a single working injection can be
reused against every customer of the same agent. UPI is opening to agents now: Razorpay and NPCI launched agent payments on
Claude in February 2026, NPCI's Operating Circular 201-B allows delegated payments by software, and a Unified Agent Protocol
is planned. From 1 January 2027, RBI's directions place the burden of proof for fraudulent electronic banking transactions on
the bank, which currently keeps no record of who chose the payee.

\section{Our approach and findings}
\textbf{Approach.} The gate examines where the payee came from rather than what the agent's inputs say. For each payment the
agent proposes, it traces the payee identifier to its source, using only records the bank already holds, and assigns one of
five classes (Figure~\ref{fig:flow}): \textbf{P}, the customer's own signed message; \textbf{K}, a payee the customer saved
earlier; \textbf{R}, the registered handle of a biller; \textbf{D}, a document the customer pointed the agent to; and
\textbf{W}, anything else. Payments to P, K and R are allowed, D requires the customer's confirmation, and W is blocked.
Later rules can only make a decision stricter. A bill that claims to come from a registered biller must carry that biller's
handle. The amount must match what the customer stated or the biller's bill fetch, and paying a saved payee more than last
time, or soon after its details changed, requires confirmation. Each decision is written to a hash-chained, HMAC-signed
receipt. The gate uses no model and makes no network call; it is 424 lines of Python and takes a median of 74\,\textmu s
per decision. We considered and rejected three alternatives. A second language model
judging each payment can be injected through the same text. A known-payee allowlist blocks every first payment to a new
payee (41.7\% of legitimate AgentDojo payments) and is bypassed by adding the payee first. A tighter spending limit still
admits small redirected payments; 80.7\% of real attacks would fit under one.

%%FIG_FLOW%%

\textbf{Threat model and resilience under exposure.} We assume the agent is fully compromised. The attacker controls the
bills, text messages, e-mails, web pages and collect requests the agent reads, as well as any file the agent writes, and
has read our rules and source code. Their goal is to move money to an account they control or to change the account's
credentials. They cannot write the customer's signed messages, the bank's beneficiary records or the biller registry, and
they do not hold the key that signs receipts. Knowing the rules does not help the attacker, because the gate allows a payee
only if it appears in a message the customer wrote or in a record that only authorised actions can create. There is no
classifier to evade and no score threshold to stay below. Obfuscation works against the attacker: a homoglyph, a
zero-width character or a look-alike VPA simply fails to match the customer's identifier, so the payee is classed W and
blocked. A 5,000-case fuzz test checks this fail-closed property, and a mutation probe confirms that each of fifteen
deliberately planted defects is caught. The only secret is the receipt key. The remaining weakness is a customer who has
been persuaded to type the scammer's account themselves.

\textbf{How we evaluated it.} We used six sources of evidence. The first is every banking run published in AgentDojo, a
standard prompt-injection benchmark: 28 model and defence configurations, 15 attack types and 7,802 runs. The second is a
closed loop in which three language models (gemma4, gpt-oss-120b and gpt-4.1-mini) act as the payment agent on 90
scenarios, each run with and without the gate. The third places the gate in front of an Indian payment gateway's
open-source payout server, built unmodified, against a local capture proxy, the provider's real sandbox API, a second provider's real API and a third's UPI sandbox (collect, Autopay, QR). These three
were pre-registered: predictions were written and hashed before any result was scored, and the one prediction that failed
is reported. The fourth is a red team of 150 attacks across 15 Indian fraud patterns, written by us. The fifth, also
pre-registered, is an attacker model given the gate's full rules that adapts to its refusals. The sixth is a set of
unit, property and mutation tests. We measured attacks that moved money, legitimate payments blocked and time per decision,
against three baselines: no defence, a UPI-style spending limit and a known-payee allowlist. All data was synthetic.

\textbf{Key results} (Figure~\ref{fig:results}; details in Appendix~\ref{app:evidence}).
\begin{itemize}
\item \textbf{AgentDojo.} Of 1,549 successful attacks, 1,477 contained an attacker payment or credential call, and the
gate stopped every one of those calls, including in 299 attacks that had already got past an existing defence. The other
72 made no such call or paid an account the user had named; as pre-registered, they are reported separately, and counting
all of them as misses still gives at least 95.4\%. A spending limit let 83.6\% through. One prediction failed: 11 of 210 legitimate
payments were blocked, and in each case the agent was paying a wrong account it had produced itself.
\item \textbf{Real models as the agent.} Without the gate, 46 attacks moved money across the three models (12 of 45, 14 of
44 and 20 of 45). With the gate, none did, and no legitimate task was lost.
\item \textbf{Real payout server and three real payment providers.} Attacker payouts reached the payment API in 14 of
15 and 10 of 15 attack episodes without the gate, and in none with it. Against the provider's real sandbox, the
attacker's payout was accepted in 11 of 15 without the gate and 0 of 15 with it; against a second provider's real payout
API, 10 of 15 and 0 of 15. Legitimate payouts were unchanged in both. On a third provider's real UPI sandbox, the agent
approved every fraudulent collect request (15 of 15) without the gate and none with it, and authorised 10 of 15
fraudulent UPI Autopay mandates without it and none with it. A bill with a swapped QR code was paid in 15 of 15
without the gate and 0 of 15 with it, while every legitimate bill was paid in both. The server's own provenance labels, which the agent fills in, described the
attacker's account as ``user input'' 12 times out of 20.
\item \textbf{An attacker who knows every rule.} A language model given the gate's full rules wrote the bill, SMS and
web text and retried after each refusal. Against a customer who reads confirmations it never succeeded (0 of 45). Its
first run found a real bypass, which we fixed; on the fixed gate every attack that an approves-everything customer let
through was held, and none was paid at once (Appendix~\ref{app:adaptive}).
\end{itemize}

%%FIG_RESULTS%%

\textbf{What's new here.} Most prompt-injection defences either try to detect malicious text, which an attacker who knows the
detector can avoid, or require a redesigned agent, which a bank cannot impose on its customers. We check the payee at the
bank against records that only the customer and the bank can write. Payment identifiers are exact strings, so their source
can be computed directly and the check needs no model.

\textbf{Who could use this, and how.} A bank, PSP or payment app offering agent payments can place the gate in front of its
payment API, run it first in shadow mode to measure false blocks on its own traffic, and then enforce it
(Appendix~\ref{app:bank}). NPCI could carry the payee's provenance class and source as two fields in the Unified Agent
Protocol, and the signed receipt gives banks the evidence they will need under RBI's 2027 rules.

\textbf{Limitations and open problems.} Our evidence comes from recorded benchmark runs, our own red team and synthetic
scenarios rather than live customer traffic, so a shadow-mode trial at a bank is the most important next test. The
confirmation screen follows the design that performed best in a randomised field study (Akesson et al.\ 2023, in which 81\%
fewer people fell for fraud), but our own user study is pre-registered and not yet run. Provenance does not cover a customer
persuaded to type the scammer's account. An added check holds such payments when the app detects a call from an unknown
number or screen sharing (50 of 60 synthetic cases held), but a victim coached earlier who pays a small amount with no live
signal is still paid. Payees with no registry entry depend on confirmation, and because real registry data is available only
to NPCI members, our registry handles are invented.

\section{Reflection}
\textbf{Difficulty faced.} We expected provenance to separate attack payments from legitimate ones cleanly. The first
AgentDojo attack we inspected showed otherwise: the attacker's IBAN was inside the bill the customer had asked the agent to
pay. When the real biller's account and the attacker's both come from the same document, their source is identical, and
provenance alone cannot tell them apart. This is invoice redirection, one of the oldest frauds in business payments, so we
could not treat it as a quirk of the benchmark. None of the obvious fixes worked. Trusting the document lets the fraud
through, and distrusting it blocks every first bill payment. Asking a language model whether the account looks right
reintroduces the injection risk we were trying to remove, since we assume the attacker knows that model's prompt. Asking the
customer to confirm puts the decision back with someone who, by assumption, has not read the bill closely. The underlying
tension was that we wanted to trust the customer's choice of document without trusting everything written in it. Working
through it changed how we framed the problem. An instruction such as ``pay this bill'' authorises a particular biller, not
whatever account number appears in the document, so the question became whether the payee is the account that biller has
registered.

\textbf{How we overcame it.}
\begin{itemize}
\item \textbf{Gave named documents their own class.} Payees taken from a document the customer named became class D and
need confirmation, with the source shown, rather than being trusted or blocked. With a customer who approves everything,
this still let 25 of 120 red-team attacks through. The cost was one extra confirmation on a first bill.
\item \textbf{Checked documents against the biller registry.} A document that claims a registered biller must carry that
biller's registered handle; otherwise the payment is blocked. Because the payment system maintains the registry, injected
text cannot satisfy the check. This reduced 25 to 5, and every legitimate payment still completed.
\item \textbf{Tested against other teams' failures.} Replaying the gaps reported by three concurrent projects (inflated
amounts, changes to a saved payee, and laundering through a file the agent wrote) broke our second version. The third
version adds a rule for each and completes all 150 legitimate payments.
\item \textbf{Made a mistaken confirmation recoverable.} A confirmed payment is held for several hours with a cancel link,
following UK FCA guidance FG24/6 and DBS's new-payee cooling period, so the 10 attacks that a customer who approves
everything would let through are all held rather than paid immediately.
\item \textbf{Still open.} Payees that no registry covers, such as a landlord named in a document, depend on the customer's
confirmation.
\end{itemize}

